\chapter{法源与核对范围}

教材为柏浪涛《2027精讲卷·刑法》。本册章、节名称对应教材的讲次及节名。
涉及规则版本时，以行为时间、规范效力及适用规则为输入。以下为正文和扩展规则主要使用的官方来源。
全书是可以逐步检查的纸面形式推导。涉及合理怀疑、情节严重、期待可能性和客观归责等判断时，仍须补入题设或依法认定的事实与评价。
有限表格只核验其中声明的条件结构。

\begin{enumerate}
 \item 《中华人民共和国刑法》，国家税务总局法规库整合文本。
 \url{https://fgk.chinatax.gov.cn/zcfgk/c100009/c5212248/content.html}
 \item 《中华人民共和国刑法修正案（十二）》及施行时间。
 \url{https://www.spp.gov.cn/spp/zdgz/202312/t20231229_638849.shtml}
 \item 《中华人民共和国刑事诉讼法》，参见第55条。
 \url{https://www.samr.gov.cn/zw/zfxxgk/fdzdgknr/bgt/art/2025/art_48b3b0c3f954448e996602a71fc6f878.html}
 \item 全国人大常委会关于《刑法》第30条的解释。
 \url{https://jjjcz.mee.gov.cn/djfg/gjflfg/flgfxwj/201404/t20140401_444600.html}
 \item 两高一部《关于依法适用正当防卫制度的指导意见》，法发〔2020〕31号。
 \url{https://www.court.gov.cn/zixun/xiangqing/251611.html}
 \item 最高检、公安部关于低龄未成年人严重暴力犯罪核准追诉案件的规定。
 \url{https://www.ah.jcy.gov.cn/jwgk/fgsd/202608/t20260803_7845099.shtml}
 \item 最高法关于审理未成年人刑事案件具体应用法律若干问题的解释。
 \url{https://www.court.gov.cn/fabu/xiangqing/1040.html}
 \item 最高法《关于处理自首和立功若干具体问题的意见》。
 \url{https://www.court.gov.cn/fabu/xiangqing/2484.html}
 \item 最高法《关于办理减刑、假释案件具体应用法律的规定》，法释〔2016〕23号。
 \url{https://www.court.gov.cn/fabu/xiangqing/30621.html}
 \item 两高罪名补充规定（八）。
 \url{https://www.court.gov.cn/fabu/xiangqing/424452.html}
 \item 两高办理强奸、猥亵未成年人刑事案件的解释，法释〔2023〕3号。
 \url{https://www.court.gov.cn/zixun/xiangqing/401002.html}
 \item 四部门《关于办理性侵害未成年人刑事案件的意见》。
 \url{https://www.court.gov.cn/shenpan/xiangqing/401012.html}
 \item 两高办理盗窃刑事案件的解释，法释〔2013〕8号。
 \url{https://www.court.gov.cn/fabu/xiangqing/5228.html}
 \item 两高办理诈骗刑事案件的解释，法释〔2011〕7号。
 \url{https://www.court.gov.cn/zixun/xiangqing/2472.html}
 \item 最高法关于审理抢劫刑事案件适用法律若干问题的指导意见。
 \url{https://www.court.gov.cn/fabu/xiangqing/37412.html}
 \item 两高办理洗钱刑事案件的解释，法释〔2024〕10号。
 \url{https://www.court.gov.cn/zixun/xiangqing/440681.html}
 \item 两高办理危害税收征管刑事案件的解释，法释〔2024〕4号。
 \url{https://www.court.gov.cn/fabu/xiangqing/428482.html}
 \item 两高办理侵犯知识产权刑事案件的解释。
 \url{https://www.court.gov.cn/zixun/xiangqing/463291.html}
 \item 两高办理掩饰、隐瞒犯罪所得及其收益刑事案件的解释。
 \url{https://www.spp.gov.cn/xwfbh/wsfbt/202508/t20250825_704537.shtml}
 \item 两高环境污染刑事解释第一次和第二次修正文本。
 \url{https://www.court.gov.cn/shenpan/xiangqing/496041.html}
 \url{https://www.court.gov.cn/fabu/xiangqing/508011.html}
 \item 两高办理贪污贿赂刑事案件的解释（二），法释〔2026〕6号。
 \url{https://www.court.gov.cn/shenpan/xiangqing/497181.html}
 \item 两高办理贪污贿赂刑事案件的解释，法释〔2016〕9号。
 \url{https://www.court.gov.cn/zixun/xiangqing/19612.html}
 \item 全国人大常委会关于《刑法》第384条第一款的解释。
 \url{https://wb.flk.npc.gov.cn/flfg/PDF/eff5a520f389458cab702ea91b2dad24.pdf}
 \item 两高办理渎职刑事案件的解释（一）。
 \url{https://www.court.gov.cn/shenpan/xiangqing/5788.html}
 \item 两高办理危害民航飞行安全刑事案件的解释。
 \url{https://www.spp.gov.cn/xwfbh/wsfbt/202604/t20260408_725512.shtml}
\end{enumerate}
